\section{Semantics}
\label{sec:semantics}

We now build the structures against which soundness is proved: a
\emph{semantics of the source} --- proof-irrelevant, classical, and
syntactic in the sense that individuals are $\lbox$ terms --- and a
first-order structure $\ML$ over the same individuals interpreting the
signature $\SigFOL$.  Throughout, $\Env$ is a fixed well-formed
definitional environment with $\Sfrag$ declarations.

\subsection{The domain}
\label{sec:semantics-domain}

\begin{definition}[Domain]
\label{def:domain}
Let $\Dom$ be the set of closed $\lbox$ terms modulo $\lconv$.  We write
$\cls{e}$ for the class of $e$.  $\Dom$ is non-empty ($\cls{\bx} \in
\Dom$) and carries an application operation $d \cdot d' \defeq
\cls{e\,e'}$ for representatives $e \in d$, $e' \in d'$, well defined
since $\lconv$ is a congruence.
\end{definition}

$\Dom$ deliberately contains \emph{junk}: classes such as
$\cls{\csO\,\csO}$ (a constructor applied to a constructor) or
$\cls{\mathsf{fix}\,f.f}$ that are the value of no source term of any
type.  By \cref{cor:discrimination}, distinct constructor forms lie in
distinct classes, and junk classes like $\cls{\csO\,\csO}$ lie outside
every constructor form.  Junk is what makes the guard analysis of
\cref{sec:casesplit} non-trivial, and keeping it in the model --- rather
than quotienting it away --- is what lets the model arbitrate which
axioms need guards.

\begin{definition}[Valuations and term denotation]
\label{def:valuation}
A \emph{valuation} $\rho$ assigns to finitely many term variables elements
of $\Dom$ and to finitely many predicate variables $q$ (of kind
$(\sigma_1, \dots, \sigma_n) \to \Prop$) functions $\Dom^n \to \{0,1\}$.
Proof variables are not assigned anything.  For a $\CICm$ term $t$ whose
free term variables are assigned by $\rho$, the \emph{denotation} is
\[
  \sem{t}_\rho \defeq \cls{\erase{t}^\rho} \in \Dom,
\]
where $\erase{t}^\rho$ substitutes a representative of $\rho(x)$ for each
free variable $x$ of $\erase{t}$.  This is well defined: free variables
of $\erase{t}$ are free term variables of $t$ (proof and predicate
variables do not survive erasure by \cref{def:erasure}), and the class is
independent of the choice of representatives because $\lconv$ is a
congruence.  For a $\lbox$ term $e$ we likewise write $\cls{e^\rho}$.
\end{definition}

\subsection{Interpretation of types, propositions and predicates}
\label{sec:semantics-interp}

\begin{definition}[Interpretations]
\label{def:interp}
By recursion on the position of declarations in $\Env$ and, within each
declaration, on the structure of phrases, we define for every set type
$\sigma$, proposition $\varphi$ and predicate expression $P$ over a prefix
of $\Env$, and every valuation $\rho$ covering their free variables:
\begin{itemize}[leftmargin=2em]
\item a set $\semr{\sigma}{\rho} \subseteq \Dom$,
\item a truth value $\pvr{\varphi}{\rho} \in \{0, 1\}$,
\item a function $\kvr{P}{\rho} : \Dom^n \to \{0,1\}$ for $P$ of kind
  $(\sigma_1, \dots, \sigma_n) \to \Prop$,
\end{itemize}
as follows.

\emph{Datatypes.}  For a datatype $I$ with constructors $\wcns{c}_i :
(\sigma_{i1}, \dots, \sigma_{ir_i})$, define stages $\stage{I}{m}
\subseteq \Dom$ by $\stage{I}{0} = \emptyset$ and
\begin{multline*}
  \stage{I}{m+1} = \bigl\{\, \cls{\wcns{c}_i(e_1, \dots, e_{r_i})} \;\bigm|\;
    i \leq \text{number of constructors},\\
    \cls{e_j} \in \semr{\sigma_{ij}}{} \text{ for non-recursive } j,\quad
    \cls{e_j} \in \stage{I}{m} \text{ for recursive } j \,\bigr\},
\end{multline*}
which is well defined because the non-recursive $\sigma_{ij}$ are closed
types over the environment prefix preceding $I$.  The stages are
increasing in $m$ (induction on $m$), and we set $\semr{I}{} =
\bigcup_{m} \stage{I}{m}$ and define the \emph{rank}
$\rank{I}{d}$ of $d \in \semr{I}{}$ as the least $m$ with $d \in
\stage{I}{m}$.

\emph{Set types.}
\begin{align*}
  \semr{\Pi x{:}\sigma.\tau}{\rho} &=
    \{\, d \in \Dom \mid \forall e \in \semr{\sigma}{\rho}.\;
      d \cdot e \in \semr{\tau}{\rho[x \mapsto e]} \,\}\\
  \semr{\varphi \to \tau}{\rho} &=
    \{\, d \in \Dom \mid \pvr{\varphi}{\rho} = 1 \implies d \in
      \semr{\tau}{\rho} \,\}\\
  \semr{\refty{x{:}\sigma}{\varphi}}{\rho} &=
    \{\, d \in \semr{\sigma}{\rho} \mid \pvr{\varphi}{\rho[x \mapsto d]}
      = 1 \,\}
\end{align*}
The clause for $\varphi \to \tau$ has no application: it mirrors the
erasure of the proof argument (\cref{def:erasure}).

\emph{Inductive predicates.}  For $J : (\sigma_1, \dots, \sigma_k) \to
\Prop$ with clauses $\wcns{k}_j : \forall \tup{x}{:}\tup{\tau}_j.\
\varphi_{j1} \to \dots \to \varphi_{j\ell_j} \to J(\tup{u}_j)$, define
stages $\pv{J}^0 = \emptyset \subseteq \Dom^k$ and
\begin{multline*}
  \pv{J}^{m+1} = \Bigl\{\, \bigl(\sem{u_{j1}}_{\rho_0}, \dots,
    \sem{u_{jk}}_{\rho_0}\bigr) \;\Bigm|\;
    j,\ \rho_0 = [\tup{x} \mapsto \tup{d}\,],\
    d_i \in \semr{\tau_{ji}}{\rho_0},\\
    \pvr{\varphi}{\rho_0} = 1 \text{ for non-rec.\ premises},\
    \tup{v}\text{-tuples} \in \pv{J}^m \text{ for rec.\ premises}
  \,\Bigr\},
\end{multline*}
where a recursive premise $J(\tup{v})$ contributes the condition
$(\sem{v_1}_{\rho_0}, \dots, \sem{v_k}_{\rho_0}) \in \pv{J}^m$.
(Premises and index types mention only $J$-free material and earlier
declarations, plus the telescope variables $\tup{x}$, so the right-hand
side is well defined.)  Set $\pv{J} = \bigcup_m \pv{J}^m$, with ranks
defined as for datatypes.

\emph{Propositions.}
\begin{align*}
  \pvr{\top}{\rho} &= 1, \qquad \pvr{\bot}{\rho} = 0\\
  \pvr{\varphi \wedge \psi}{\rho} &= \min(\pvr{\varphi}{\rho},
    \pvr{\psi}{\rho}), \qquad
  \pvr{\varphi \vee \psi}{\rho} = \max(\pvr{\varphi}{\rho},
    \pvr{\psi}{\rho})\\
  \pvr{\varphi \to \psi}{\rho} &= \max(1 - \pvr{\varphi}{\rho},
    \pvr{\psi}{\rho})\\
  \pvr{\forall x{:}\sigma.\varphi}{\rho} &=
    \min_{d \in \semr{\sigma}{\rho}} \pvr{\varphi}{\rho[x \mapsto d]},
  \qquad
  \pvr{\exists x{:}\sigma.\varphi}{\rho} =
    \max_{d \in \semr{\sigma}{\rho}} \pvr{\varphi}{\rho[x \mapsto d]}\\
  \pvr{\forall q{:}K.\varphi}{\rho} &=
    \min_{Q : \Dom^n \to \{0,1\}} \pvr{\varphi}{\rho[q \mapsto Q]}\\
  \pvr{a =_\sigma b}{\rho} &= 1 \text{ iff } \sem{a}_\rho = \sem{b}_\rho\\
  \pvr{J(\tup{a})}{\rho} &= 1 \text{ iff }
    (\sem{a_1}_\rho, \dots, \sem{a_k}_\rho) \in \pv{J}\\
  \pvr{q(\tup{a})}{\rho} &= \rho(q)(\sem{a_1}_\rho, \dots,
    \sem{a_n}_\rho)
\end{align*}
(with $\min \emptyset = 1$, $\max \emptyset = 0$).

\emph{Predicate expressions.}  $\kvr{q}{\rho} = \rho(q)$ and
$\kvr{\lambda(\tup{x}{:}\tup{\sigma}).\varphi}{\rho}(\tup{d}) =
\pvr{\varphi}{\rho[\tup{x} \mapsto \tup{d}\,]}$.
\end{definition}

The recursion is well founded: the datatype and predicate stages refer
only to earlier declarations and to the (structurally smaller) phrases of
their own declaration; all other clauses recurse on proper sub-phrases.
The semantics is classical and \emph{proof-irrelevant} by construction:
truth values do not depend on proofs, and $\sem{\cdot}$ ignores proof
subterms.  Note also that the interpretation of equality is untyped ---
$\pv{a =_\sigma b}$ compares classes in $\Dom$ and does not mention
$\sigma$ --- which matches the translation clause $\F{a =_\sigma b} =
\compf(\erase{a}) \feq \compf(\erase{b})$.

\begin{lemma}[Elementary properties]
\label{lem:sem-basic}
\begin{enumerate}[leftmargin=2em]
\item $\semr{\sigma}{\rho}$, $\pvr{\varphi}{\rho}$, $\kvr{P}{\rho}$
  depend only on the restriction of $\rho$ to the free variables of the
  phrase.
\item (Weakening) They are unchanged by extending $\rho$ on fresh
  variables.
\item Every $d \in \semr{I}{}$ has a finite rank, and
  $d \in \stage{I}{m}$ for all $m > \rank{I}{d}$; similarly for
  $\pv{J}$.
\end{enumerate}
\end{lemma}

\begin{proof}
(1), (2): routine induction on the phrase.  (3): by definition of the
union and monotonicity of stages.
\end{proof}

\begin{lemma}[Semantic substitution]
\label{lem:sem-subst}
Let $a$ be a term, $\pi$ a proof and $P$ a predicate expression, with
$\rho$ covering the relevant free variables.  Then:
\begin{enumerate}[leftmargin=2em]
\item $\sem{\subst{t}{a}{x}}_\rho = \sem{t}_{\rho[x \mapsto
  \sem{a}_\rho]}$, and the analogous equations for the substituted
  type $\subst{\sigma}{a}{x}$, proposition $\subst{\varphi}{a}{x}$ and
  predicate expression $\subst{P}{a}{x}$.
\item Substitution of proofs is invisible:
  $\sem{\subst{t}{\pi}{p}}_\rho = \sem{t}_\rho$,
  $\semr{\subst{\sigma}{\pi}{p}}{\rho} = \semr{\sigma}{\rho}$,
  $\pvr{\subst{\varphi}{\pi}{p}}{\rho} = \pvr{\varphi}{\rho}$.
\item $\pvr{\subst{\varphi}{P}{q}}{\rho} = \pvr{\varphi}{\rho[q \mapsto
  \kvr{P}{\rho}]}$, and analogously for types and terms (on which the
  substitution acts only through annotations and is invisible).
\end{enumerate}
\end{lemma}

\begin{proof}
(1) For terms: $\erase{\subst{t}{a}{x}} =
\subst{\erase{t}}{\erase{a}}{x}$ by \cref{lem:erase-subst}, and
substituting a representative of $\sem{a}_\rho$ for $x$ yields a
representative of the same class.  For types, propositions and
predicates: induction on the phrase, with the term case as base; the
binder cases use (2) of \cref{lem:sem-basic}; the predicate-application
case $\subst{(q(\tup{a}))}{P}{q} = P \cdot \tup{a}$ in (3) unfolds
$\kvr{P}{\rho}$ by definition.  (2) By \cref{lem:erase-subst}, erasure
discards the substituted proof, and proofs do not occur in the semantic
clauses.  (3) As in (1).
\end{proof}

\begin{lemma}[Conversion invariance]
\label{lem:conv-inv}
If $t \conv t'$ then $\sem{t}_\rho = \sem{t'}_\rho$.  If $\sigma \conv
\sigma'$ (both well formed) then $\semr{\sigma}{\rho} =
\semr{\sigma'}{\rho}$; if $\varphi \conv \varphi'$ then
$\pvr{\varphi}{\rho} = \pvr{\varphi'}{\rho}$.
\end{lemma}

\begin{proof}
For terms: by \cref{lem:erase-sim}, $t \conv t'$ implies $\erase{t}
\lconv \erase{t'}$, hence $\erase{t}^\rho \lconv \erase{t'}^\rho$
($\lconv$ is closed under substitution of representatives), i.e.\
$\sem{t}_\rho = \sem{t'}_\rho$.

For types and propositions, first consider a single step $\sigma \red
\sigma'$.  By \cref{lem:rigid} (and its proof), the step occurs inside a
term position (an equation side, a predicate-atom argument, an index) or
inside a component type or proposition.  Proceed by induction on the
structure of $\sigma$: in the first case the claim follows from the term
statement, since the semantic clauses depend on embedded terms only
through $\sem{\cdot}_\rho$; in the second, from the induction hypothesis
applied to the component (for binders, at every extension $\rho[x
\mapsto d]$, using that the same step relates the instantiated
components).  The general claim follows by induction on the length of the
conversion path.
\end{proof}

\subsection{The first-order structure \texorpdfstring{$\ML$}{M}}
\label{sec:semantics-model}

\begin{definition}[The structure $\ML$]
\label{def:model}
$\ML$ is the first-order $\SigFOL$-structure with universe $\Dom$ and:
\begin{itemize}[leftmargin=2em]
\item $\feq$ interpreted as identity of classes; $\fapp$ as $\cdot$;
  $\prfc$ as $\cls{\bx}$;
\item $\Ihat{I}^{\ML} = \semr{I}{}$ for each datatype $I$, and
  $\Ihat{J}^{\ML} = \pv{J}$ for each inductive predicate $J$;
\item $\Ihat{\wcns{c}}^{\ML}(d_1, \dots, d_r) = \cls{\wcns{c}(e_1, \dots,
  e_r)}$ for representatives $e_i \in d_i$;
\item for every symbol $f$ introduced by compilation with closure arity
  $s$, extension arity $m$ and witness $e_f$ (a $\lbox$ term with free
  variables $\tup{w} = w_1, \dots, w_s$; see
  \cref{def:compC,def:compile}):
  \[
    f^{\ML}(d_1, \dots, d_s, d'_1, \dots, d'_m) =
      \cls{\,\subst{e_f}{\tup{e}}{\tup{w}}\; e'_1 \cdots e'_m\,}
  \]
  for representatives $e_i \in d_i$, $e'_j \in d'_j$.  In particular
  $\Ihat{c}^{\ML}(\tup{d}) = \cls{c\; \tup{e}}$ (the $\lbox$ constant $c$
  applied) and $(\Ihat{c}^0)^{\ML} = \cls{c}$.
\end{itemize}
All interpretations are well defined on classes by congruence of
$\lconv$.  We write $\den{u}{\rho}$ for the denotation in $\ML$ of a
first-order term $u$ under an assignment $\rho$ of its variables, and
$\ML, \rho \vDash \phi$ for satisfaction of a first-order formula, in the
standard (classical, Tarskian) sense.
\end{definition}

Since source term variables double as first-order variables, a valuation
$\rho$ (\cref{def:valuation}) restricted to term variables serves as a
first-order assignment, and we use the same letter for both.

\subsection{Adequacy of the term translation and of compilation}
\label{sec:semantics-adequacy}

\begin{definition}[Coherent template maps]
\label{def:coherent}
Let $\theta$ be a template map as in \cref{sec:translation-compile}, and
$\sigma_\theta$ a substitution assigning to each $g \in \dom{\theta}$ a
closed $\lbox$ term $\sigma_\theta(g)$.  Say $(\theta, \sigma_\theta)$ is
\emph{coherent} for an assignment $\rho$ if for every $g$ with $\theta(g)
= (f, \tup{w}, m)$ and all $d_1, \dots, d_m \in \Dom$ with
representatives $e_i$,
\[
  f^{\ML}\bigl(\den{\tup{w}}{\rho},\, \tup{d}\bigr) =
  \cls{\,\sigma_\theta(g)\;e_1 \cdots e_m\,}.
\]
\end{definition}

\begin{lemma}[Adequacy of $\compf$]
\label{lem:comp-adequacy}
Let $e$ be a $\lbox$ term whose free variables are first-order variables
assigned by $\rho$ or bound by $\theta$, and let $(\theta,
\sigma_\theta)$ be coherent for $\rho$.  Then
\[
  \den{\compf_\theta(e)}{\rho} = \cls{\,e^{\rho, \sigma_\theta}\,},
\]
where $e^{\rho,\sigma_\theta}$ substitutes representatives of $\rho(x)$
for the first-order variables and $\sigma_\theta(g)$ for the
$\theta$-bound variables.
\end{lemma}

\begin{proof}
By induction on $e$, following the clauses of \cref{def:compC}.
Variables, $\bx$ and constructors are immediate.  For a spine $h\,e_1
\cdots e_j$:
\begin{itemize}[leftmargin=2em]
\item $h = g$ $\theta$-bound: by coherence,
  $f^{\ML}(\den{\tup{w}}{\rho}, \den{\compf_\theta(e_1)}{\rho}, \dots) =
  \cls{\sigma_\theta(g)\, e_1^{\rho,\sigma_\theta} \cdots}$ using the
  induction hypothesis on the $e_i$; the remaining arguments are appended
  by $\fapp$, interpreted as $\cdot$.
\item $h = c$: $\Ihat{c}^{\ML}$ applies the $\lbox$ constant $c$ to the
  argument classes (\cref{def:model}), which is exactly
  $\cls{(c\,e_1\cdots e_j)^{\rho,\sigma_\theta}}$, with the pointer case
  identical via $(\Ihat{c}^0)^{\ML} = \cls{c}$ and $\fapp$.
\item $h$ a variable: immediate.
\item $h$ complex: by \cref{def:model} the head's value translation
  denotes $\cls{h^{\rho,\sigma_\theta}}$ --- for a $\lambda$-lift,
  $f^{\ML}(\den{\tup{w}}{\rho}) = \cls{e_f^{\,\rho,\sigma_\theta}}$ with
  $e_f = h$, and similarly for case lifts ($g^{\ML}(\den{\compf(s)}{},
  \dots) = \cls{(\lambda z.\mathsf{case}(z;\tup{br}))^{\dots}\,
  s^{\dots}}$, which $\beta$-reduces to
  $\cls{\mathsf{case}(s;\tup{br})^{\dots}}$, using the induction
  hypothesis for $s$) and fix lifts (via the pointer's witness);
  the spine arguments are then appended by $\cdot$.
\end{itemize}
Note the case for complex heads uses only the \emph{interpretation} of
the lifted symbols, not the axioms emitted for them.
\end{proof}

\begin{lemma}[Soundness of compilation]
\label{lem:compile-sound}
Consider any invocation of the compilation procedure on the environment
(or on a term embedded in the goal), and let
$\mathrm{B}_\theta(L; e)$ range over the body-compilation calls it
performs.  Then:
\begin{enumerate}[leftmargin=2em]
\item (\emph{Invariant}) For every such call there is a substitution
  $\sigma_\theta$ such that for every assignment $\rho$ of the variables
  of $L$, the pair $(\theta, \sigma_\theta)$ is coherent for $\rho$ and
  $\den{L}{\rho} = \cls{e^{\rho, \sigma_\theta}}$.
\item Every axiom emitted by the procedure holds in $\ML$.
\end{enumerate}
\end{lemma}

\begin{proof}
By induction on the tree of procedure calls.  We first check the
invariant at the roots, then that each rule preserves it and that each
emission is valid given it.

\emph{Roots.}  For a definition $c \defeq t : \tau$ with $\erase{t} =
\lambda \tup{z}.e_0$ (no fix): $\den{\Ihat{c}(\tup{z})}{\rho} =
\cls{c\,\tup{e}} = \cls{\erase{t}\,\tup{e}}$ (a $\delta$-step) $=
\cls{e_0^{\rho}}$ (by $|\tup{z}|$ $\beta$-steps), with $\theta =
\emptyset$.  For $\erase{t} = \mathsf{fix}\,g_0.\lambda\tup{z}.e_0$: set
$\sigma_\theta(g_0) = \erase{t}$; coherence holds because
$\Ihat{c}^{\ML}(\tup{d}) = \cls{c\,\tup{e}} =
\cls{\erase{t}\,\tup{e}}$ by a $\delta$-step; and $\den{\Ihat{c}(\tup{z})}
{\rho} = \cls{c\,\tup{e}} = \cls{\erase{t}\,\tup{e}} =
\cls{(\lambda\tup{z}.e_0)[\erase{t}/g_0]\,\tup{e}} =
\cls{e_0^{\rho,\sigma_\theta}}$ by a $\mu$-step and $\beta$-steps.  For a
$\lambda$-lift of $h = \lambda x.e'$ (\cref{def:compC}), the call is
$\mathrm{B}_\theta(f(\tup{w}); h)$ and $\den{f(\tup{w})}{\rho} =
\cls{h^{\rho,\sigma_\theta}}$ holds by the definition of $f^{\ML}$
(witness $h$); for a case lift, $\den{g(z, \tup{w}')}{\rho} =
\cls{(\lambda z'.\mathsf{case}(z';\tup{br}))^{\dots}\, z^\rho}
\lconv \cls{\mathsf{case}(z;\tup{br})^{\rho,\sigma_\theta}}$ by a
$\beta$-step; for a fix lift $h = \mathsf{fix}\,g_0.\lambda
\tup{z}.e_0$, extend $\sigma_\theta$ with $\sigma_\theta(g_0) =
h^{\rho,\sigma_\theta}|_{\tup{w}}$ --- the closed instance of $h$ at the
representatives of $\den{\tup{w}}{\rho}$; coherence for $g_0$ is the
definition of $f^{\ML}$ via witness $h$, and the invariant for
$\mathrm{B}_{\theta'}(f(\tup{w},\tup{z}); e_0)$ follows by one $\mu$-step
and $\beta$-steps as in the definition-root case.

\emph{Rule (1)} ($e = \lambda x.e'$): $\den{L \fapp x}{\rho} =
\den{L}{\rho} \cdot \rho(x) = \cls{e^{\dots}\,x^\rho} \lconv
\cls{e'^{\,\rho,\sigma_\theta}}$ by a $\beta$-step.

\emph{Rule (2)} (pattern split): an assignment $\rho'$ of the refined
left-hand side $\subst{L}{\Ihat{\wcns{c}}_i(\tup{z}_i)}{x}$ induces the
assignment $\rho = \rho'[x \mapsto \cls{\wcns{c}_i(\tup{e})}]$ of $L$;
then
\[
  \den{\subst{L}{\Ihat{\wcns{c}}_i(\tup{z}_i)}{x}}{\rho'} =
  \den{L}{\rho} = \cls{\mathsf{case}(x;\dots)^{\rho,\sigma_\theta}},
\]
and a head $\iota$-step reduces the representative to
$\cls{(\subst{b_i}{\wcns{c}_i(\tup{z}_i)}{x})^{\rho',\sigma_\theta}}$.

\emph{Rule (3)} (flatten): the emitted axiom $L \feq
g(\compf_\theta(s), \tup{w}')$ is valid: by the invariant and
\cref{lem:comp-adequacy}, both sides denote
$\cls{\mathsf{case}(s;\tup{br})^{\rho,\sigma_\theta}}$ (for the right
side, $g^{\ML}$ applies its witness $\lambda z.\mathsf{case}(z;\tup{br})$
to $\den{\compf_\theta(s)}{\rho} = \cls{s^{\rho,\sigma_\theta}}$ and one
$\beta$-step concludes).  The recursive call's invariant is the root
computation for case lifts above.

\emph{Rule (4)} (inner fix): the emitted axiom $L \feq f^0(\tup{w})$ is
valid since both sides denote $\cls{e^{\rho,\sigma_\theta}}$ (the right
side by the definition of $(f^0)^{\ML}$ with witness $e$); the recursive
compilation is covered by the fix-lift root case.

\emph{Rule (5)}: the emitted axiom $L \feq \compf_\theta(e)$ is valid by
the invariant and \cref{lem:comp-adequacy}.

\emph{Pointer axioms}: $\forall \tup{z}.\ \Ihat{c}^0 \fapp z_1 \fapp
\cdots \fapp z_{a_c} \feq \Ihat{c}(\tup{z})$ holds because both sides
denote $\cls{c\,\tup{e}}$; similarly for fix-lift pointers.
\end{proof}

\subsection{Coincidence: the translation reflects the semantics}
\label{sec:semantics-coincidence}

\begin{lemma}[Coincidence]
\label{lem:coincidence}
Let $\sigma$ be a shallow set type and $\varphi$ a shallow proposition
(over $\Env$ and a context of term variables), and let $\rho$ be a
valuation of their free variables.  Then:
\begin{enumerate}[leftmargin=2em]
\item for every first-order term $u$ over the variables of $\rho$:\quad
  $\ML, \rho \vDash \guard{\sigma}{u} \iff \den{u}{\rho} \in
  \semr{\sigma}{\rho}$;
\item $\ML, \rho \vDash \F{\varphi} \iff \pvr{\varphi}{\rho} = 1$.
\end{enumerate}
\end{lemma}

\begin{proof}
Mutual induction on the structure of $\sigma$ and $\varphi$, following
\cref{def:guard-F}.
\begin{itemize}[leftmargin=2em]
\item $\sigma = I$: $\guard{I}{u} = \Ihat{I}(u)$ and $\Ihat{I}^{\ML} =
  \semr{I}{}$ by \cref{def:model}.
\item $\sigma = \Pi x{:}\sigma'.\tau$: $\ML, \rho \vDash \forall x.\
  \guard{\sigma'}{x} \to \guard{\tau}{u \fapp x}$ iff for all $d \in
  \Dom$ with $d \in \semr{\sigma'}{\rho}$ (induction hypothesis (1) at the
  variable $x$) we have $\den{u}{\rho} \cdot d \in \semr{\tau}{\rho[x
  \mapsto d]}$ (induction hypothesis (1) for $\tau$, at the assignment
  $\rho[x \mapsto d]$ and the term $u \fapp x$), which is the definition
  of $\den{u}{\rho} \in \semr{\Pi x{:}\sigma'.\tau}{\rho}$.
\item $\sigma = \varphi' \to \tau$ and $\sigma =
  \refty{x{:}\sigma'}{\varphi'}$: as above, using induction hypothesis
  (2) for $\varphi'$; the refinement case uses the standard first-order
  substitution lemma, $\ML, \rho \vDash
  \subst{\F{\varphi'}}{u}{x} \iff \ML, \rho[x \mapsto \den{u}{\rho}]
  \vDash \F{\varphi'}$.
\item Connectives, $\top$, $\bot$: the classical satisfaction clauses
  match the truth-value clauses of \cref{def:interp} literally.
\item $\F{\forall x{:}\sigma'.\varphi'} = \forall x.\guard{\sigma'}{x}
  \to \F{\varphi'}$: satisfaction quantifies over all $d \in \Dom$ and
  relativises by the guard, which by induction hypothesis (1) carves out
  exactly $\semr{\sigma'}{\rho}$; this matches
  $\min_{d \in \semr{\sigma'}{\rho}}$.  Dually for $\exists$.
\item Equations: $\ML, \rho \vDash \compf(\erase{a}) \feq
  \compf(\erase{b})$ iff $\den{\compf(\erase{a})}{\rho} =
  \den{\compf(\erase{b})}{\rho}$ iff $\cls{\erase{a}^\rho} =
  \cls{\erase{b}^\rho}$ (by \cref{lem:comp-adequacy} with $\theta =
  \emptyset$) iff $\sem{a}_\rho = \sem{b}_\rho$, which is
  $\pvr{a =_\sigma b}{\rho} = 1$.
\item $J$-atoms: as for equations, using $\Ihat{J}^{\ML} = \pv{J}$.
\end{itemize}
Shallowness is used exactly here: predicate quantifiers and
predicate-variable atoms have no translation, and all remaining formers
are covered.
\end{proof}
